\Needspace{12\baselineskip}
\begin{center}
\textit{Einteilung der Klassen in Ordnungen.}
\end{center}

\Needspace{8\baselineskip}
\Article{226}
Wir nennen eine Form $(a,b,c)$ \textit{primitiv}, wenn die Zahlen $a,b,c$ keinen gemeinsamen nichttrivialen Teiler haben; andernfalls heißt sie \textit{abgeleitet}. Genauer: Setzen wir den größten gemeinsamen Teiler der Zahlen $a,b,c$ $=m$, so ist die Form $(a,b,c)$ \textit{aus der primitiven Form} $\left(\frac{a}{m},\frac{b}{m},\frac{c}{m}\right)$ \textit{abgeleitet}. Aus dieser Definition ist sofort klar, dass alle Formen, deren Determinante durch keine Quadratzahl (außer $1$) teilbar ist, notwendig primitiv sind. Ferner folgt aus Artikel 161: Enthält eine gegebene Klasse von Formen mit der Determinante $D$ eine primitive Form, so sind alle Formen dieser Klasse primitiv; in diesem Fall
\SourceLeaf{232}
heißt die Klasse selbst \textit{primitiv}. Ist außerdem eine Form $F$ mit der Determinante $D$ aus einer primitiven Form $f$ mit der Determinante $\frac{D}{mm}$ abgeleitet, und sind $F$, $f$ jeweils in den Klassen $K$, $k$ enthalten, so sind offenbar alle Formen der Klasse $K$ aus der primitiven Klasse $k$ abgeleitet. Daher nennen wir in diesem Fall die Klasse $K$ selbst \textit{aus der primitiven Klasse $k$ abgeleitet}.

Ist $(a,b,c)$ eine primitive Form, sind aber $a$, $c$ nicht beide gerade (\textit{d.h.}, wenn entweder beide ungerade sind oder wenigstens eine davon ungerade ist), so sieht man leicht, dass nicht nur $a$, $b$, $c$, sondern auch $a$, $2b$, $c$ keinen gemeinsamen nichttrivialen Teiler haben. Dann heißt die Form $(a,b,c)$ \textit{eigentlich primitiv} oder einfach eine \textit{eigentliche Form}. Ist dagegen $(a,b,c)$ eine primitive Form und sind die Zahlen $a$, $c$ beide gerade, so haben die Zahlen $a$, $2b$, $c$ offenbar den gemeinsamen Teiler $2$ (der zugleich der größte ist). Dann heißt $(a,b,c)$ eine \textit{uneigentlich primitive Form} oder einfach eine \textit{uneigentliche Form}\HistoricalNote{*)}{Wir haben die Ausdrücke \textit{eigentlich} und \textit{uneigentlich} hier gewählt, weil uns keine passenderen einfielen. Wir weisen darauf hin, damit niemand zwischen dieser Bedeutung und der seit Artikel 157 verwendeten einen verborgenen Zusammenhang sucht: Es gibt keinen. Eine Mehrdeutigkeit ist hier jedenfalls nicht zu befürchten.}. In diesem Fall muss $b$ ungerade sein (sonst wäre $(a,b,c)$ keine primitive Form). Daher gilt $bb\equiv1\ (\mathrm{mod.}\ 4)$, und da $ac$ durch $4$ teilbar ist, gilt für die Determinante $bb-ac\equiv1\ (\mathrm{mod.}\ 4)$. Uneigentliche Formen gibt es also nur für eine Determinante der Form $4n+1$, wenn sie positiv ist, oder der Form $-(4n+3)$, wenn sie negativ ist. --- Aus Artikel 161 ist außerdem klar: Enthält eine gegebene Klasse eine eigentlich primitive Form, so sind alle Formen dieser Klasse eigentlich primitiv. Dagegen besteht eine Klasse, die eine uneigentlich primitive Form enthält, ausschließlich aus uneigentlich primitiven Formen. Daher heißt die Klasse selbst im ersten Fall \textit{eigentlich primitiv} oder einfach \textit{eigentlich}, im zweiten \textit{uneigentlich primitiv} oder \textit{uneigentlich}. So sind \textit{zum Beispiel} unter den positiven Klassen von Formen mit der Determinante $-235$ sechs eigentlich, nämlich die mit den Repräsentanten $(1,0,235)$, $(4,1,59)$, $(4,-1,59)$, $(5,0,47)$, $(13,5,20)$, $(13,-5,20)$, und ebenso viele unter den negativen Klassen; jeweils zwei sind dagegen uneigentlich. --- Die Klassen von Formen mit der Determinante $79$ (einer Zahl der Form $4n+3$) sind sämtlich eigentlich.

Ist eine Form $(a,b,c)$ abgeleitet, und zwar aus der primitiven Form $\left(\frac{a}{m},\frac{b}{m},\frac{c}{m}\right)$, so kann diese eigentlich oder uneigentlich primitiv sein. Im ersten Fall ist $m$ auch der größte gemeinsame Teiler der Zahlen $a$, $2b$, $c$; im zweiten ist deren größter gemeinsamer Teiler $2m$. Damit wird die Unterscheidung zwischen einer \textit{aus einer eigentlich primitiven Form abgeleiteten Form} und einer \textit{aus einer uneigentlich primitiven Form abgeleiteten Form} verständlich; ebenso (da sich nach Artikel 161 alle Formen derselben Klasse in dieser Hinsicht gleich verhalten)
\SourceLeaf{233}
die Unterscheidung zwischen einer \textit{aus einer eigentlich primitiven Klasse abgeleiteten Klasse} und einer \textit{aus einer uneigentlich primitiven Klasse abgeleiteten Klasse}.

Mit diesen Unterscheidungen haben wir eine erste Grundlage, auf der wir die Einteilung aller Klassen von Formen mit gegebener Determinante in verschiedene \textit{Ordnungen} aufbauen können. Zwei Klassen mit den repräsentierenden Formen $(a,b,c)$, $(a',b',c')$ rechnen wir zur \textit{selben Ordnung}, wenn \textit{sowohl} die Zahlen $a$, $b$, $c$ denselben größten gemeinsamen Teiler haben wie $a'$, $b'$, $c'$ \textit{als auch} $a$, $2b$, $c$ denselben wie $a'$, $2b'$, $c'$. Ist eine oder sind beide dieser Bedingungen nicht erfüllt, so gehören die Klassen zu \textit{verschiedenen Ordnungen}. Daraus ist sofort klar, dass alle eigentlich primitiven Klassen eine Ordnung bilden und alle uneigentlich primitiven eine andere. Ist $mm$ eine Quadratzahl, die die Determinante $D$ teilt, so bilden die aus eigentlich primitiven Klassen mit der Determinante $\frac{D}{mm}$ abgeleiteten Klassen eine eigene Ordnung, die aus uneigentlich primitiven Klassen mit der Determinante $\frac{D}{mm}$ abgeleiteten eine andere usw. Ist $D$ durch keine Quadratzahl (außer $1$) teilbar, so gibt es keine Ordnungen abgeleiteter Klassen. Es gibt dann entweder nur eine Ordnung (wenn $D\equiv2$ oder $3$ modulo $4$), nämlich die Ordnung der eigentlich primitiven Klassen, oder zwei (wenn $D\equiv1\ (\mathrm{mod.}\ 4)$), nämlich die Ordnung der eigentlich primitiven und die der uneigentlich primitiven Klassen. Nach den Grundsätzen der Kombinatorik lässt sich leicht folgende allgemeine Regel aufstellen: Es sei $D=D'2^{2\mu}a^{2\alpha}b^{2\beta}c^{2\gamma}\ldots.$, wobei $D'$ keinen quadratischen Faktor enthält und $a$, $b$, $c$ usw. verschiedene ungerade Primzahlen sind (auf diese Form lässt sich jede Zahl bringen, indem man $\mu=0$ setzt, wenn $D$ nicht durch $4$ teilbar ist, und $\alpha$, $\beta$, $\gamma$ usw. sämtlich $=0$ setzt oder, was dasselbe ergibt, die Faktoren $a^{2\alpha}$, $b^{2\beta}$, $c^{2\gamma}$ usw. weglässt, wenn $D$ durch keine ungerade Quadratzahl teilbar ist). Dann gibt es entweder
\[
(\mu+1)(\alpha+1)(\beta+1)(\gamma+1)\ldots
\]
Ordnungen, nämlich wenn $D'\equiv2$ oder $3\ (\mathrm{mod.}\ 4)$, oder
\[
(\mu+2)(\alpha+1)(\beta+1)(\gamma+1)\ldots
\]
Ordnungen, wenn $D'\equiv1\ (\mathrm{mod.}\ 4)$. Den Beweis dieser Regel lassen wir jedoch weg, da er weder schwierig noch hier besonders notwendig ist.

\textit{Beispiel} 1. Für $D=45=5.3^2$ gibt es sechs Klassen mit den Repräsentanten $(1,0,-45)$, $(-1.\,0,45)$, $(2,1,-22)$, $(-2,1,22)$, $(3.\,0,-15)$, $(6,3,-6)$. Sie werden in vier Ordnungen eingeteilt: Ordnung I umfasst die beiden eigentlichen Klassen mit den Repräsentanten $(1,0,-45)$, $(-1,0,45)$; Ordnung II enthält die beiden uneigentlichen Klassen\SourceLeaf{234} mit den Repräsentanten $(2,1,-22)$, $(-2,1,22)$; Ordnung III enthält eine aus einer eigentlichen Klasse der Determinante $5$ abgeleitete Klasse, nämlich die mit dem Repräsentanten $(3,0,-15)$; Ordnung IV besteht aus einer aus einer uneigentlichen Klasse der Determinante $5$ abgeleiteten Klasse, nämlich der mit dem Repräsentanten $(6,3,-6)$.

\textit{Beispiel} 2. Die positiven Klassen mit der Determinante $-99=-11.3^2$ werden in vier Ordnungen eingeteilt. Ordnung I umfasst die folgenden eigentlich primitiven Klassen\HistoricalNote{*)}{Der Kürze wegen verwenden wir die repräsentierenden Formen anstelle der Klassen, die sie vertreten.}: $(1,0,99)$, $(4,1,25)$. $(4,-1,25)$, $(5,1,20)$, $(5,-1,20)$, $(9,0,11)$; Ordnung II enthält die uneigentlichen Klassen $(2,1,50)$, $(10,1,10)$; Ordnung III die aus eigentlichen Klassen der Determinante $-11$ abgeleiteten Klassen $(3,0,33)$, $(9,3,12)$, $(9,-3,12)$; Ordnung IV die einzige aus einer uneigentlichen Klasse der Determinante $-11$ abgeleitete Klasse $(6,3,18)$. --- Die negativen Klassen dieser Determinante lassen sich genauso in Ordnungen einteilen.

Wir bemerken, dass \textit{entgegengesetzte Klassen stets zur selben Ordnung gehören}; der Grund für diesen Satz ist leicht zu erkennen.

\Needspace{8\baselineskip}
\Article{227}
Von diesen verschiedenen Ordnungen verdient besonders die Ordnung der eigentlich primitiven Klassen die größte Aufmerksamkeit. Denn jede abgeleitete Klasse stammt aus einer bestimmten primitiven Klasse (mit kleinerer Determinante); aus deren Betrachtung ergibt sich das auf die erste Klasse Bezügliche meist von selbst. Unten werden wir zeigen, dass jede uneigentlich primitive Klasse ähnlich gewissermaßen entweder einer einzigen eigentlich primitiven Klasse oder dreien (mit derselben Determinante) zugeordnet ist. Bei negativen Determinanten können wir außerdem die negativen Klassen übergehen, da jeder von ihnen stets eine bestimmte positive Klasse entspricht. Um das Wesen der eigentlich primitiven Klassen tiefer zu verstehen, erläutern wir zunächst einen bestimmten wesentlichen Unterschied, nach dem sich die ganze Ordnung der eigentlichen Klassen in mehrere \textit{Geschlechter} einteilen lässt. Da wir diesen sehr wichtigen Gegenstand bisher noch nicht berührt haben, müssen wir ihn von Anfang an aufnehmen.

\Needspace{12\baselineskip}
\begin{center}
\textit{Einteilung der Ordnungen in Geschlechter.}
\end{center}

\Needspace{8\baselineskip}
\Article{228}
\textsc{Satz.} \textit{Durch jede eigentlich primitive Form $F$ lassen sich unendlich viele Zahlen darstellen, die durch eine beliebige gegebene Primzahl $p$ nicht teilbar sind.}

\textit{Beweis.} Ist die Form $F=axx+2bxy+cyy$, so kann $p$ offenbar nicht alle drei Zahlen $a$, $2b$, $c$ zugleich teilen. Ist nun $a$ nicht durch $p$ teilbar\SourceLeaf{235}, so ist klar: Nimmt man für $x$ eine beliebige nicht durch $p$ teilbare Zahl und für $y$ eine durch $p$ teilbare Zahl, so ist der Wert der Form $F$ nicht durch $p$ teilbar. Ist $c$ nicht durch $p$ teilbar, so erreicht man dasselbe, indem man $x$ einen teilbaren und $y$ einen nicht teilbaren Wert gibt. Sind schließlich sowohl $a$ als auch $c$ durch $p$ teilbar, also $2b$ nicht teilbar, so nimmt die Form $F$ einen nicht durch $p$ teilbaren Wert an, wenn man sowohl $x$ als auch $y$ beliebige nicht durch $p$ teilbare Werte gibt. \textit{Q.E.D.}

Der Satz gilt offenbar auch für \textit{uneigentlich primitive} Formen, sofern nur nicht $p=2$ gilt.

Da mehrere solche Bedingungen gleichzeitig erfüllt sein können, dass dieselbe Zahl durch gewisse gegebene Primzahlen teilbar, durch andere nicht teilbar ist (siehe Artikel 32), sieht man leicht: Die Zahlen $x$, $y$ lassen sich auf unendlich viele Arten so bestimmen, dass die primitive Form $axx+2bxy+cyy$ einen Wert erhält, der durch beliebig viele gegebene Primzahlen nicht teilbar ist. Von diesen ist nur $2$ auszuschließen, wenn die Form uneigentlich primitiv ist. Daraus folgt diese allgemeinere Fassung des Satzes: \textit{Durch jede primitive Form lassen sich unendlich viele Zahlen darstellen, die zu einer beliebigen gegebenen Zahl teilerfremd sind (die gegebene Zahl muss ungerade sein, wenn die Form uneigentlich primitiv ist).}

\Needspace{8\baselineskip}
\Article{229}
\textsc{Satz.} \textit{Sei $F$ eine primitive Form der Determinante $D$ und $p$ eine Primzahl, die $D$ teilt. Dann stimmen die nicht durch $p$ teilbaren Zahlen, die durch die Form $F$ darstellbar sind, darin überein, dass sie entweder alle quadratische Reste von $p$ oder alle Nichtreste sind.}

\textit{Beweis.} Sei $F=(a,b,c)$; seien $m$, $m'$ zwei beliebige nicht durch $p$ teilbare Zahlen, die durch die Form $F$ darstellbar sind, nämlich
\[
m=agg+2bgh+chh,\qquad m'=ag'g'+2bg'h'+ch'h'
\]
Dann ist
\[
mm'=(agg'+b(gh'+hg')+chh')^2-D(gh'-hg')^2
\]
Also ist $mm'$ einem Quadrat modulo $D$ und damit auch modulo $p$ kongruent, \textit{d.h.}, $mm'$ ist quadratischer Rest von $p$. Daraus folgt, dass entweder beide Zahlen $m$, $m'$ quadratische Reste von $p$ oder beide Nichtreste sind. \textit{Q.E.D.}

\SourceLeaf{236}
Ebenso beweist man: Ist die Determinante $D$ durch $4$ teilbar, so sind alle durch $F$ darstellbaren ungeraden Zahlen entweder sämtlich $\equiv1$. oder sämtlich $\equiv3\ (\mathrm{mod.}\ 4)$. Denn das Produkt zweier solcher Zahlen ist in diesem Fall stets quadratischer Rest von $4$ und daher $\equiv1\ (\mathrm{mod.}\ 4)$; somit sind entweder beide $\equiv1$ oder beide $\equiv3$.

Ist schließlich $D$ durch $8$ teilbar, so ist das Produkt zweier beliebiger ungerader Zahlen, die durch $F$ darstellbar sind, quadratischer Rest von $8$ und folglich $\equiv1\ (\mathrm{mod.}\ 8)$. Daher sind in diesem Fall alle durch $F$ darstellbaren ungeraden Zahlen entweder sämtlich $\equiv1$ oder sämtlich $\equiv3$ oder sämtlich $\equiv5$ oder sämtlich $\equiv7\ (\mathrm{mod.}\ 8)$.

So sind \textit{zum Beispiel}, da die Form $(10,3,17)$ die Zahl $10$ darstellt, die ein Nichtrest von $7$ ist, alle nicht durch $7$ teilbaren Zahlen, die jene Form darstellt, Nichtreste von $7$. --- Da $-3$ durch die Form $(-3,1,49)$ darstellbar und modulo $4$ $\equiv1$ ist, verhalten sich alle durch diese Form darstellbaren ungeraden Zahlen ebenso.

Wäre es für unseren gegenwärtigen Zweck nötig, so könnten wir leicht beweisen, dass die durch die Form $F$ darstellbaren Zahlen zu keiner Primzahl, die $D$ nicht teilt, eine solche feste Beziehung haben. Vielmehr lassen sich sowohl Reste als auch Nichtreste jeder Primzahl, die $D$ nicht teilt, durch die Form $F$ darstellen. Bei den Zahlen $4$ und $8$ gilt dagegen auch in anderen Fällen etwas Ähnliches, das wir nicht übergehen können.

I. \textit{Ist die Determinante $D$ der primitiven Form $F$ $\equiv3\ (\mathit{mod.}\ 4)$, so sind alle durch die Form $F$ darstellbaren ungeraden Zahlen entweder sämtlich $\equiv1$ oder sämtlich $\equiv3\ (\mathit{mod.}\ 4)$.} Sind nämlich $m$, $m'$ zwei durch $F$ darstellbare Zahlen, so lässt sich ihr Produkt $mm'$ wie oben auf die Form $pp-Dqq$ bringen. Sind beide Zahlen $m$, $m'$ ungerade, so muss von den Zahlen $p$, $q$ eine gerade und die andere ungerade sein; daher ist von den Quadraten $pp$, $qq$ eines $\equiv0$ und das andere $\equiv1\ (\mathrm{mod.}\ 4)$. Daraus folgt leicht, dass $pp-Dqq$ sicher $\equiv1\ (\mathrm{mod.}\ 4)$ ist, also entweder beide Zahlen $m$, $m'$ $\equiv1$ sind oder beide $\equiv3\ (\mathrm{mod.}\ 4)$. So lassen sich \textit{zum Beispiel} durch die Form $(10,3,17)$ keine anderen ungeraden Zahlen als solche der Form $4n+1$ darstellen.

II. \textit{Ist die Determinante $D$ der primitiven Form $F$ $\equiv2\ (\mathit{mod.}\ 8)$, so sind alle durch die Form $F$ darstellbaren ungeraden Zahlen entweder teils $\equiv1$, teils $\equiv7$, oder teils $\equiv3$, teils $\equiv5\ (\mathit{mod.}\ 8)$.} Seien nämlich $m$, $m'$ zwei ungerade Zahlen,
\SourceLeaf{237}
die durch $F$ darstellbar sind; ihr Produkt $mm'$ lässt sich also auf die Form $pp-Dqq$ bringen. Sind daher beide Zahlen $m$, $m'$ ungerade, so muss $p$ ungerade sein (weil $D$ gerade ist), also gilt $pp\equiv1\ (\mathrm{mod.}\ 8)$. Dagegen ist $qq$ entweder $\equiv0$ oder $\equiv1$ oder $\equiv4$, und daher $Dqq$ entweder $\equiv0$ oder $\equiv2$. Somit ist $mm'=pp-Dqq$ entweder $\equiv1$ oder $\equiv7\ (\mathrm{mod.}\ 8)$. Ist also $m$ entweder $\equiv1$ oder $\equiv7$, so ist auch $m'$ entweder $\equiv1$ oder $\equiv7$; ist dagegen $m$ entweder $\equiv3$ oder $\equiv5$, so ist auch $m'$ entweder $\equiv3$ oder $\equiv5$. \textit{Zum Beispiel} sind alle durch die Form $(3,1,5)$ darstellbaren ungeraden Zahlen entweder $\equiv3$ oder $\equiv5\ (\mathrm{mod.}\ 8)$, und keine Zahlen der Form $8n+1$ oder $8n+7$ sind durch jene Form darstellbar.

III. \textit{Ist die Determinante $D$ der primitiven Form $F$ $\equiv6\ (\mathit{mod.}\ 8)$, so sind durch diese Form entweder nur ungerade Zahlen der beiden Arten $\equiv1$ und $\equiv3\ (\mathit{mod.}\ 8)$ oder nur solche der beiden Arten $\equiv5$ und $\equiv7\ (\mathit{mod.}\ 8)$ darstellbar.} Einen Beweis, der dem vorigen (in II) ganz ähnlich ist, kann jeder ohne Mühe ausführen. --- So sind \textit{zum Beispiel} durch die Form $(5,1,7)$ nur solche ungeraden Zahlen darstellbar, die entweder $\equiv5$ oder $\equiv7\ (\mathrm{mod.}\ 8)$ sind.

\Needspace{8\baselineskip}
\Article{230}
Alle Zahlen, die durch eine gegebene primitive Form $F$ der Determinante $D$ darstellbar sind, haben also eine feste Beziehung zu jedem Primteiler von $D$ (durch den sie selbst nicht teilbar sind). Die durch $F$ darstellbaren ungeraden Zahlen haben in gewissen Fällen auch zu den Zahlen $4$ und $8$ eine feste Beziehung: nämlich zu $4$, wenn $D$ entweder $\equiv0$ oder $\equiv3\ (\mathrm{mod.}\ 4)$ ist, und zu $8$, wenn $D$ entweder $\equiv0$ oder $\equiv2$ oder $\equiv6\ (\mathrm{mod.}\ 8)$ ist\HistoricalNote{*)}{Für durch $8$ teilbare Determinanten kann die Beziehung zur Zahl $4$ weggelassen werden, da sie in diesem Fall schon in der Beziehung zu $8$ enthalten ist.}. Eine solche Beziehung zu jeder dieser Zahlen nennen wir einen \textit{Charakter} oder \textit{Einzelcharakter} der Form $F$ und drücken sie so aus: Können nur quadratische Reste der Primzahl $p$ durch die Form $F$ dargestellt werden, so geben wir ihr den Charakter $Rp$, im entgegengesetzten Fall den Charakter $Np$. Ebenso schreiben wir $1,4$, wenn die Form $F$ keine anderen ungeraden Zahlen darstellen kann als solche, die $\equiv1\ (\mathrm{mod.}\ 4)$ sind. Daraus ist sofort klar, welche Charaktere die Zeichen $3,4$; $1,8$; $3,8$; $5,8$; $7,8$ ausdrücken. Schließlich geben wir den Formen, die nur solche ungeraden Zahlen darstellen können, die modulo
\SourceLeaf{238}
$8$ entweder $\equiv1$ oder $\equiv7$ sind, den Charakter $1\ \mathit{und}\ 7,8$. Daraus ergibt sich von selbst die Bedeutung der Charaktere $3\ \mathit{und}\ 5,8$; $1\ \mathit{und}\ 3,8$; $5\ \mathit{und}\ 7,8$.

Die einzelnen Charaktere einer gegebenen primitiven Form $(a,b,c)$ mit der Determinante $D$ lassen sich stets aus wenigstens einer der Zahlen $a$, $c$ erkennen (die offenbar beide durch jene Form darstellbar sind). Denn ist $p$ ein Primteiler von $D$, so ist sicher wenigstens eine der Zahlen $a$, $c$ nicht durch $p$ teilbar. Wären beide durch $p$ teilbar, so würde $p$ auch $bb(=D+ac)$ teilen. und damit auch $b$; \textit{d.h.}, die Form $(a,b,c)$ wäre nicht primitiv. Ebenso ist in den Fällen, in denen die Form $(a,b,c)$ zur Zahl $4$ oder $8$ eine feste Beziehung hat, sicher mindestens eine der Zahlen $a$, $c$ ungerade. Daraus lässt sich also jene Beziehung erkennen. So ergibt sich \textit{zum Beispiel} für die Form $(7,0,23)$ der Charakter bezüglich der Zahl $23$ aus der Zahl $7$ als $N23$; der Charakter derselben Form bezüglich $7$ ergibt sich aus der Zahl $23$, nämlich $R7$; schließlich lässt sich der Charakter dieser Form bezüglich $4$, nämlich $3,4$, sowohl aus $7$ als auch aus $23$ erkennen.

Da alle Zahlen, die durch eine Form $F$ der Klasse $K$ darstellbar sind, auch durch jede andere Form dieser Klasse dargestellt werden können, kommen die einzelnen Charaktere der Form $F$ offenbar auch allen übrigen Formen dieser Klasse zu. Wir können sie daher als Charaktere der ganzen Klasse betrachten. Die einzelnen Charaktere einer beliebigen gegebenen primitiven Klasse lassen sich also aus ihrer repräsentierenden Form erkennen. Entgegengesetzte Klassen haben stets sämtliche Charaktere gemeinsam.
